
% Lecture Notes: Variant Calling and Downstream Analysis of NGS Data

\section{Introduction to Downstream Applications of NGS}

Once sequencing reads have been mapped and aligned to a reference genome, the data can be used to extract a variety of biological insights. While the primary focus of this lecture is the identification of genomic variants, several other powerful applications rely on the foundations of read mapping.

\subsection{Alternative Sequencing Applications}

By modifying the experimental protocol before sequencing, we can target specific biological phenomena rather than just mapping the entire genomic sequence.

\begin{important}
  \textbf{ChIP-seq (Chromatin Immunoprecipitation Sequencing)}: A method used to analyze protein interactions with DNA. Instead of sequencing the entire genome, only the DNA fragments to which a specific protein (e.g., a transcription factor like TP53) has bound are isolated and sequenced.
\end{important}

In ChIP-seq, mapped reads concentrate strongly around specific genomic regions where the protein physically interacted with the DNA, while the rest of the genome exhibits only low-level "random noise". Researchers use statistical \textit{peak calling} algorithms to identify these high-density regions and map transcription factor binding sites.

\begin{important}
  \textbf{RNA-seq (RNA Sequencing)}: A technique that sequences transcribed RNA molecules (converted to complementary DNA, or cDNA) to measure gene expression levels and analyze transcript structure.
\end{important}

Because RNA-seq samples mature RNA, the mapped reads align almost exclusively to exonic (protein-coding) regions. The analysis provides two main types of information:
\begin{itemize}
    \item \textbf{Expression levels:} Genes with higher read coverage correspond to higher expression levels in the cell.
    \item \textbf{Alternative splicing:} Reads can span splice junctions, starting in one exon and ending in another, effectively skipping the intronic regions that were spliced out.
\end{itemize}

% [IMAGE: A "Sashimi" plot showing RNA-seq read coverage over exonic regions. Arcs connect exons to indicate splice-junction spanning reads. A comparison between a control sample and a patient sample illustrates aberrant splicing caused by a mutation, leading to skipped exons.]

\begin{important}
  \textbf{Hi-C}: A chromosome conformation capture method that identifies three-dimensional chromatin interactions by fusing interacting genomic regions and sequencing the resulting artificial junctions.
\end{important}

Hi-C produces paired-end reads where one read maps to a specific location (e.g., Chromosome 2) and its pair maps to a completely different location (e.g., Chromosome 4), indicating that these two regions were physically adjacent in the cell nucleus. The output is typically visualized as a two-dimensional interaction map.

\section{Genomic Variations and Variant Calling}

The main application of Whole Genome Sequencing (WGS) is identifying differences between the sequenced genome of an individual and the standard reference genome. 

\begin{important}
  \textbf{Variant Calling}: The computational process of identifying variable sites (mutations) and their genomic locations relative to a reference genome.
\end{important}

\begin{important}
  \textbf{Genotype Calling}: Determining the exact genotype for an individual at each identified variable site (e.g., homozygous reference, heterozygous variant, or homozygous variant).
\end{important}

We distinguish between two fundamentally different types of mutations based on their origin:
\begin{itemize}
    \item \textbf{Germline Mutations:} Present in the germ cells (egg or sperm) and therefore inherited and present in \textit{every} cell of the resulting organism's body. \sta These are responsible for hereditary disorders, such as Fibrodysplasia Ossificans Progressiva (a severe disease causing soft tissue to ossify, caused by a single point mutation in the ACVR1 gene).
    \item \textbf{Somatic Mutations:} Acquired after conception in specific somatic (non-germline) cells. \sta Somatic mutations are the primary drivers of cancer (e.g., ERBB2 mutations in glioblastoma). They are identified by comparing a tumor sample to a "matched normal" control sample (e.g., blood) from the same patient to filter out harmless germline variants.
\end{itemize}

\subsection{Types of Genomic Variations}

Genomic variations range from single-base changes to massive structural rearrangements.

\begin{important}
  \textbf{SNP (Single-Nucleotide Polymorphism)}: A single-base variant that is relatively common, being present in at least 1\% of a given population. Pronounced "snip".
\end{important}

\begin{important}
  \textbf{SNV (Single-Nucleotide Variant)}: A single-base variant identified in an individual genome, regardless of its frequency in the population.
\end{important}

\sta In medical genomics, researchers are often interested in rare SNVs that are not established SNPs, as common SNPs are unlikely to be the sole cause of a rare, severe genetic disease. A typical human genome contains about 3–4 million SNVs.

\begin{important}
  \textbf{Indels}: Small insertions or deletions of one or a few nucleotides.
\end{important}

\textbf{Example:} Frameshift Mutation
Indels in protein-coding regions can cause disastrous "frameshifts." If a sequence is read in triplets (e.g., \texttt{CAT CAT CAT}), inserting a single 'A' shifts the entire reading frame (\texttt{cAt aCA Tca...}), changing all subsequent amino acids and likely truncating the protein. This illustrates why even small indels can be highly destructive.

\begin{important}
  \textbf{Structural Variants (SVs)}: Large-scale genomic alterations, typically defined as being 1 kbp or larger, including deletions, duplications, inversions, and translocations.
\end{important}

SVs are categorized as:
\begin{itemize}
    \item \textbf{Balanced}: No net change in the amount of DNA (e.g., inversions, translocations).
    \item \textbf{Unbalanced / Copy Number Variants (CNVs)}: The amount of DNA is increased or decreased (e.g., deletions, duplications). A human genome typically contains hundreds or thousands of SVs, with an average size of 250,000 nucleotides—much larger than an average gene.
\end{itemize}

\section{Evaluating Sequencing Confidence}

To reliably call variants, we must quantify our confidence in the sequencing data using two critical metrics: mapping quality and base quality.

\begin{important}
  \textbf{Mapping Quality}: The confidence with which a given read maps to a unique, specific genomic location. A low mapping quality indicates the read maps ambiguously to multiple highly repetitive regions.
\end{important}

\begin{important}
  \textbf{Base Quality}: The confidence with which an individual nucleotide was sequenced by the machine (e.g., how unambiguous the optical signal was for that specific base).
\end{important}

Both qualities are typically expressed as \textbf{PHRED quality scores}, computed using a simple logarithmic transformation:
\begin{equation}
    Q_{PHRED} = -10 \log_{10} p
\end{equation}
where $p$ is the probability that the base call or read mapping is \textit{wrong}.

\textbf{Example:} PHRED Score Interpretation
A PHRED score of $Q=30$ implies $p = 10^{-3} = 0.001$. This means there is a 1 in 1000 probability of error, corresponding to 99.9\% accuracy.

\subsection{Variant Allele Frequency (VAF)}

\begin{important}
  \textbf{Variant Allele Frequency (VAF)}: The fraction of mapped reads at a specific genomic position that support the alternate (mutated) base rather than the reference base.
\end{important}

For germline mutations, VAF clusters around expected discrete values:
\begin{itemize}
    \item \textbf{Heterozygous variant:} $\approx 50\%$ of reads (since one of two chromosomal copies is affected).
    \item \textbf{Homozygous variant:} $\approx 100\%$ of reads (since both chromosomal copies are affected).
\end{itemize}

\sta In cancer (somatic mutations), VAF analysis is much more complex due to \textit{tumor purity} (normal cells mixed with tumor cells) and \textit{tumor clonality} (sub-populations of tumor cells carrying different mutations). For instance, a mutation present in only 5\% of reads might represent a small but critical clone responsible for therapy resistance.

\section{Algorithms for Single Nucleotide Variant Calling}

\subsection{Alignment Pileups}
To evaluate variants, mapped reads are visually and computationally stacked into a "pileup" format. 

\begin{minted}{text}
21 31587791 T 24 .,.,.,.,.,.,.,.,.,.,.,.,^],
21 31587792 G 25 .,.,.,.,.,.,.,.,.,.,.,^],
21 31587793 A 26 .,.,g.Gg,.G..GG,G.gG.,g,^],
21 31587794 G 27 .,.,.,.,.,.,.,.,.,.,.,.,^],
\end{minted}

This output from `samtools` represents a sequence alignment column-by-column. 
\begin{enumerate}
    \item Each line shows the chromosome, 1-based coordinate, and the reference base (e.g., `A` at position 31587793).
    \item The number indicates read coverage (26 reads cover position 31587793).
    \item Dots `.` and commas `,` represent a match to the reference base on the forward and reverse strands, respectively.
    \item Letters (`ACGTN` or `acgtn`) indicate mismatches. At position 31587793, 10 reads favor `G` (alternate) and 16 favor `A` (reference), strongly suggesting a heterozygous variant.
\end{enumerate}

\subsection{A Naive Frequency-Based Heuristic}

Early NGS studies used strict quality cutoffs (e.g., filtering out bases with $Q < 20$) followed by a straightforward frequency threshold:
\begin{itemize}
    \item If VAF $< 20\% \implies$ Homozygous reference.
    \item If $20\% \le$ VAF $\le 80\% \implies$ Heterozygous variant.
    \item If VAF $> 80\% \implies$ Homozygous variant.
\end{itemize}

\sta \textbf{Limitations:} This naive approach fails severely at low sequencing depths. If a position has only 6 reads, observing 1 variant read yields a 17\% VAF (called "reference"), while 2 variant reads yield a 33\% VAF (called "heterozygous"). This is an insignificant statistical difference resulting from random sampling. Furthermore, throwing away data strictly below a quality threshold destroys useful probabilistic information.

\subsection{Probabilistic Approach: The MAQ Algorithm}

Modern algorithms like MAQ (Mapping and Assembly with Qualities) use Bayesian inference to model the uncertainty of both the sequencing and mapping processes.

\begin{important}
  \textbf{Bayes' Theorem}: Relates the conditional probability of a model given the data, to the likelihood of the data given the model:
  \[ P(M \mid D) = \frac{P(D \mid M) P(M)}{P(D)} \]
\end{important}

MAQ aims to find the \textbf{Maximum A Posteriori (MAP)} estimate—the genotype $g \in \{ \langle a,a \rangle, \langle a,b \rangle, \langle b,b \rangle \}$ that maximizes the posterior probability:
\begin{equation}
    \hat{g} = \underset{g}{\operatorname{argmax}} p(g \mid D) = \underset{g}{\operatorname{argmax}} p(D \mid g) p(g)
\end{equation}

\subsubsection{Integrating Mapping and Base Quality}
Before computing genotype probabilities, MAQ adjusts the base quality $Q_i$ of a nucleotide based on the mapping quality $Q_z$ of the read it belongs to:
\begin{equation}
    Q_i = \min(Q_i, Q_z)
\end{equation}
\textit{Rationale:} If a base was sequenced perfectly ($Q_i$ is high) but the read might map to the wrong chromosome ($Q_z$ is low), we must temper our confidence in that base. Conversely, a confident mapping does not save a poorly sequenced base.

\subsubsection{Genotype Prior Probabilities ($P(g)$)}
MAQ sets a prior probability $r$ for observing a heterozygous genotype. To reduce false positives, $r$ is low ($r = 0.001$) for novel mutations, but much higher ($r = 0.2$) for known SNPs.
\[ P(\langle a, a \rangle) = (1 - r) / 2 \]
\[ P(\langle b, b \rangle) = (1 - r) / 2 \]
\[ P(\langle a, b \rangle) = r \]

\subsubsection{Likelihood Calculation ($P(D \mid g)$)}
For a column of $n$ mapped reads with $k$ reference bases ($a$) and $n-k$ alternate bases ($b$):
\begin{itemize}
    \item \textbf{Heterozygous case $\langle a,b \rangle$}: Modeled as a standard binomial process where drawing either allele has a $p=0.5$ probability.
    \[ P(D \mid \langle a,b \rangle) = \binom{n}{k} \frac{1}{2^n} \]
    \item \textbf{Homozygous reference $\langle a,a \rangle$}: Every observed alternate base $b$ must be a sequencing/mapping error. Let $\alpha_{n,n-k}$ be the probability of observing $n-k$ errors in $n$ bases.
    \item \textbf{Homozygous alternate $\langle b,b \rangle$}: Every observed reference base $a$ must be an error. Let $\alpha_{n,k}$ be the probability of $k$ errors.
\end{itemize}

MAQ calculates the product of the likelihood and the prior for all three genotypes, normalizes them so they sum to 1, and selects the genotype with the highest posterior probability.

\section{Somatic Variant Calling in Cancer}

Calling somatic mutations requires distinguishing true tumor-specific variants from background germline variants. Algorithms like \textbf{VarScan2} accomplish this by processing a tumor sample and a matched normal control sample in parallel.

\subsection{The VarScan2 Heuristic and Fisher's Exact Test}
If the independently called genotypes of the tumor and normal samples disagree (e.g., Normal = Homozygous Reference, Tumor = Heterozygous Variant), VarScan2 tests whether the difference in allele counts is statistically significant.

It constructs a $2 \times 2$ contingency table of read counts (Reference vs. Alternate) for both samples, and applies a one-tailed \textbf{Fisher's exact test} utilizing the hypergeometric distribution:
\begin{equation}
    P(X \ge N_{T,alt}) = \sum_{i=N_{T,alt}}^{N_T} \frac{\binom{N_T}{i} \binom{N_C}{N_{C,alt}}}{\binom{N_{tot}}{i+N_{C,alt}}}
\end{equation}
If the resulting p-value is $\le 0.05$, the variant is officially called a \textit{somatic mutation}.

\begin{important}
  \textbf{Loss of Heterozygosity (LOH)}: An event where a normal cell is heterozygous (e.g., alleles A and T), but the tumor sample shows only one allele (e.g., T). This often occurs because the genomic segment containing the 'A' allele was physically deleted in the tumor's DNA.
\end{important}

\section{Structural Variations and CNV Calling}

Structural variants, specifically Copy Number Variants (CNVs), are detected not by observing mismatched bases, but by observing the \textbf{read depth} (coverage) across regions.

% [IMAGE: Coverage plot comparing a tumor sample and a control sample. The y-axis shows log2-normalized read counts. The tumor plot shows a massive spike indicating an amplification (e.g., 60 copies of a gene) and sudden drops indicating deletions of entire chromosome arms.]

\subsection{Read Depth and the Poisson Distribution}
The genome is divided into fixed-size windows (e.g., 1 kbp). The expected number of reads mapping to a window follows a Poisson distribution:
\begin{equation}
    P(X = k) = \frac{\lambda^k e^{-\lambda}}{k!}
\end{equation}
where $\lambda = \frac{N}{G} \cdot W$ (Total Reads / Genome Size $\times$ Window Size).

If a window normally expects $\lambda$ reads, a heterozygous deletion (1 copy instead of 2) yields roughly $0.5\lambda$ reads, and a duplication (3 copies) yields $1.5\lambda$ reads. For sufficiently large $\lambda$, the Poisson distribution is approximated by a normal distribution, allowing us to compute a $Z$-score for each window:
\begin{equation}
    z = \frac{x - \mu}{\sigma}
\end{equation}
which is converted into an upper-tail or lower-tail p-value to test for abnormal increases or decreases in coverage.

\subsection{Event-Wise Testing (EWT) and Multiple Testing Correction}

To prevent calling spurious CNVs from a single noisy window, algorithms search for \textit{intervals} of $\ell$ consecutive windows that show a collective deviation. 

\begin{important}
  \textbf{Multiple Testing Correction (MTC)}: A statistical adjustment required when performing millions of simultaneous tests across the genome. Without correction, a 5\% false-positive rate ($p=0.05$) over 50 tests yields an expected error rate of $1 - (0.95)^{50} \approx 92.3\%$, guaranteeing false discoveries.
\end{important}

The most stringent MTC is the \textbf{Bonferroni correction}, which divides the target significance level ($\alpha=0.05$) by the total number of tests ($k$):
\[ p_{threshold} = \frac{\alpha}{k} \]

In Yoon et al.'s Event-Wise Testing (EWT) algorithm, an interval $\mathcal{A}$ of $\ell$ windows is called a duplication if the maximum (least significant) p-value inside the interval meets a modified, less conservative threshold:
\begin{equation}
    \max_i \{ P_i^{\text{Upper}} \mid i \in \mathcal{A} \} < \left( \frac{\ell}{L} \times \text{FPR} \right)^{\frac{1}{\ell}}
\end{equation}
Here, $L$ is the total number of windows on the chromosome, and FPR is the target false positive rate (e.g., 0.05). The term $\frac{\ell}{L} \times \text{FPR}$ effectively acts as a Bonferroni correction adjusted for the length of the event.
